\section{Teste de usabilidade com especialistas}
\label{sec:protocolo-avaliacao}\label{sec:avaliacao-especialistas}

O teste de usabilidade foi feito com um formulário respondido por um painel de especialistas, chamado aqui de formulário dos especialistas. O formulário é compartilhado com o texto do coautor: foi enviado ao mesmo grupo de respondentes, e cada trabalho analisa os itens que interessam ao seu tema. Aqui interessam os itens sobre a interface, o assistente de personalização e os fluxos do ambiente virtual de aprendizagem. O formulário foi aplicado entre 1º e 9 de novembro de 2026, recebeu oito respostas e está reproduzido no Apêndice~\ref{ap:formulario-especialistas}. As primeiras subseções descrevem o instrumento e o procedimento; as seguintes apresentam os resultados de cada bloco, já com a discussão. O resumo agregado das respostas está no Apêndice~\ref{ap:respostas}.

\subsection{Quem responde e o que é possível concluir}

A avaliação principal foi feita por um painel de especialistas, formado por professores e profissionais de computação. Neste trabalho, considera-se especialista quem atua ou já atuou com programação, como docente de disciplinas da área ou como profissional de desenvolvimento de software. Não é um estudo com os alunos, que são o público-alvo da plataforma. Pediu-se ao respondente que avaliasse se a ferramenta é adequada para um aluno em primeiro contato com programação, e não que relatasse a própria experiência de aprendizagem.

Essa escolha limita o que o trabalho pode afirmar, e por isso ela é explicada antes da descrição do instrumento. Com um painel de especialistas é possível dizer que profissionais da área acham plausível que a personalização da sintaxe reduza a barreira sintática. Não é possível dizer que ela reduz. A pergunta feita na introdução continua sem resposta empírica, e respondê-la exige medir o desempenho de estudantes, como descrito na Seção~\ref{subsec:limites-interpretacao}. A sessão com estudantes não muda esse limite, porque também colheu opinião.

Em troca, o painel dá acesso a avaliações que um aluno iniciante não teria como fazer: se o subconjunto da linguagem é adequado a uma disciplina introdutória, como a plataforma se compara às ferramentas discutidas no Capítulo~\ref{cap:trabalhos-relacionados} e quais riscos pedagógicos a proposta traz. Há ainda um motivo ligado ao tema deste trabalho. Personalizar a linguagem é uma tarefa de autoria: alguém precisa decidir quais palavras substituem quais e verificar se o vocabulário final faz sentido. Na escola, quem costuma fazer isso é o professor, que prepara o ambiente antes da aula. Perguntar a professores se o assistente é claro e se o trabalho de preparação compensa é, portanto, perguntar a quem de fato usa essa parte da plataforma.

\subsection{Composição do instrumento}

O formulário é respondido pelo próprio participante, foi distribuído pela internet e foi criado por um programa, para que sua estrutura possa ser gerada de novo sempre que necessário. A Tabela~\ref{tab:blocos-pa01} mostra os blocos na ordem em que aparecem. Três deles vêm de instrumentos já consolidados na literatura de avaliação de interfaces, e os outros nove foram escritos para este projeto. O primeiro é a \textit{System Usability Scale} (SUS), de \citeonline{brooke1996}, que mede a facilidade de uso percebida com dez afirmações e resume as respostas em um escore de 0 a 100. O segundo é a versão curta do \textit{User Experience Questionnaire} (UEQ-S), de \citeonline{schrepp2017}, que avalia a impressão geral sobre o produto por meio de pares de adjetivos opostos. O terceiro vem do \textit{Technology Acceptance Model} (TAM), de \citeonline{davis1989}, cujas afirmações tratam da utilidade percebida e da intenção de uso. O funcionamento e o cálculo de cada um estão na Subseção~\ref{subsec:escalas-escores}.

\begin{table}[htbp]
\centering
\caption{Blocos do formulário dos especialistas, procedência e número de itens.}
\label{tab:blocos-pa01}
\footnotesize
\renewcommand{\arraystretch}{1.2}
\begin{tabular}{|p{5.2cm}|p{4.6cm}|c|}
\hline
\textbf{Bloco} & \textbf{Procedência} & \textbf{Itens} \\\hline
Termo de consentimento (texto de leitura, sem item de marcação) & Elaborado para o projeto & -- \\\hline
Perfil profissional & Elaborado para o projeto & 4 \\\hline
Uso em sala (condicional) & Elaborado para o projeto & 10 \\\hline
Como a plataforma foi usada & Elaborado para o projeto & 4 \\\hline
Avaliação dos recursos & Elaborado para o projeto & 10 \\\hline
Facilidade de uso & SUS \cite{brooke1996} & 10 \\\hline
Impressão geral & UEQ-S \cite{schrepp2017} & 8 \\\hline
Potencial pedagógico & Elaborado para o projeto & 8 \\\hline
Riscos e limitações & Elaborado para o projeto & 3 \\\hline
Qualidade técnica & Elaborado para o projeto & 6 \\\hline
Utilidade e intenção de uso & TAM \cite{davis1989} & 6 \\\hline
Comentários abertos & Elaborado para o projeto & 5 \\\hline
\end{tabular}
\legend{Fonte: elaborado pelo autor.}
\end{table}

A Tabela~\ref{tab:blocos-pa01} corresponde ao formulário aplicado, cujo texto completo está no Apêndice~\ref{ap:formulario-especialistas}. Em relação à versão descrita na revisão anterior deste texto, duas afirmações mudaram de bloco e passaram a ser respondidas só por quem leciona: a que trata de alunos da mesma turma com sintaxes diferentes, antes no bloco de riscos, e a que pergunta se o subconjunto da linguagem é suficiente para uma disciplina introdutória, antes no bloco de qualidade técnica. Com isso o bloco de uso em sala passou de oito para dez itens, o de riscos ficou com três afirmações e o de qualidade técnica com cinco afirmações e uma pergunta de múltipla escolha sobre as construções que fizeram falta.

O SUS, o UEQ-S e o TAM foram usados porque têm valores de referência publicados. Com eles é possível comparar o resultado obtido com o de outros produtos já avaliados, em vez de apenas descrever a opinião dos respondentes. Os blocos próprios foram escritos porque nenhum instrumento pronto pergunta sobre o que esta plataforma propõe. Antes do bloco de avaliação dos recursos, o formulário mostra o mesmo programa em duas versões: uma na configuração padrão, com \texttt{int}, \texttt{for}, \texttt{if} e delimitadores simbólicos, e outra personalizada, com \texttt{numero\_inteiro}, \texttt{para}, \texttt{se}, operadores escritos por extenso, os delimitadores \texttt{inicio} e \texttt{fim} e nenhum ponto e vírgula. As duas versões estão transcritas no Apêndice~\ref{ap:formulario-especialistas}. Assim o respondente avalia algo concreto, e não uma ideia abstrata de personalização.

\subsection{Itens sobre a interface e a personalização}

No bloco de avaliação dos recursos, cada afirmação é respondida em uma escala de concordância de cinco pontos, de 1 (discordo totalmente) a 5 (concordo totalmente), e cada linha tem também a opção ``não usei'', o que separa quem discorda de quem não chegou a experimentar o recurso. As afirmações são:

\begin{itemize}
    \item criar ou ajustar uma linguagem no assistente foi simples;
    \item ficou claro em que etapa o usuário estava;
    \item o exemplo de código ajudou a entender o efeito de cada alteração;
    \item escrever programas com a linguagem configurada foi confortável;
    \item as mensagens de erro ajudaram a identificar o que corrigir;
    \item resolver e enviar exercícios foi simples;
    \item importar uma linguagem do acervo da comunidade foi simples.
\end{itemize}
 O bloco termina com três itens sobre a utilidade da tabela de \textit{tokens}, do código intermediário e da execução passo a passo.

Os três primeiros itens correspondem a decisões de interface descritas na Seção~\ref{sec:customizacao-dinamica}: dividir a personalização em etapas, indicar em qual delas o usuário está e mostrar uma pré-visualização do código. Esses itens verificam escolhas que foram adotadas como hipótese durante o desenvolvimento. O bloco de qualidade técnica traz mais duas afirmações sobre o assistente: se ele é compreensível para quem não tem formação em compiladores e se a validação prévia impede configurações de linguagem inválidas. Nesse bloco, a opção fora da escala é ``não sei''.

Quem declara que leciona, ou já lecionou, disciplina que envolva programação responde a um bloco extra sobre o ambiente virtual de aprendizagem, com a opção ``não se aplica'' em cada linha. Ele pergunta:

\begin{itemize}
    \item se foi simples criar uma turma e receber alunos pelo código de acesso;
    \item se foi simples criar um exercício com casos de teste;
    \item se travar uma linguagem em um exercício ajuda a planejar as aulas;
    \item se a correção automática ajuda a acompanhar a turma;
    \item se o respondente usaria a plataforma em uma disciplina introdutória de programação;
    \item se o respondente conseguiria preparar uma aula sem apoio técnico;
    \item se o trabalho de preparar a linguagem compensa o ganho em sala;
    \item se alunos de uma mesma turma usando sintaxes diferentes dificultariam o trabalho do professor;
    \item se o subconjunto da linguagem é suficiente para uma disciplina introdutória.
\end{itemize}
 A penúltima afirmação é um item de risco e é lida como tal. Uma questão de múltipla escolha pergunta em que momento de uma disciplina introdutória a plataforma seria mais útil, e uma das alternativas diz que não há espaço para ela na disciplina. Essa alternativa foi mantida de propósito, pelo mesmo motivo que justifica o bloco de riscos.

O bloco de potencial pedagógico, respondido em terceira pessoa, transforma em afirmações as premissas que sustentam a plataforma:

\begin{itemize}
    \item deixar o aluno definir as palavras-chave tende a reduzir a barreira sintática no primeiro contato;
    \item usar palavras em português no lugar de termos em inglês facilita entender o código;
    \item a personalização tende a reduzir os erros de sintaxe de iniciantes;
    \item mostrar os erros no vocabulário escolhido pelo próprio aluno facilita a correção;
    \item o conjunto de elementos configuráveis é adequado ao propósito da ferramenta.
\end{itemize}
 Os outros três itens do bloco perguntam se expor a tabela de \textit{tokens}, o código intermediário e a execução passo a passo tem valor didático para quem está aprendendo.

\subsection{Escalas usadas e cálculo dos escores}
\label{subsec:escalas-escores}

Os blocos em escala usam dois formatos de resposta. O SUS e os blocos de uso em sala, avaliação dos recursos, potencial pedagógico, riscos, qualidade técnica e utilidade e intenção de uso pedem o grau de concordância com cada afirmação em cinco pontos, em que 1 corresponde a ``discordo totalmente'', 3 ao ponto neutro e 5 a ``concordo totalmente''. Nos blocos elaborados para o projeto há também uma opção fora da escala, que não entra nas médias: ``não se aplica'' no uso em sala, ``não usei'' na avaliação dos recursos e ``não sei'' nos blocos de potencial pedagógico, de riscos e de qualidade técnica. O UEQ-S usa sete pontos entre dois adjetivos opostos. As perguntas de perfil, as de uso da plataforma, a do momento da disciplina e a das construções que fizeram falta são de múltipla escolha. O último bloco tem quatro questões abertas e um item de múltipla escolha em que o respondente diz se autoriza a citação das suas respostas sem identificação.

A facilidade de uso usa a escala de \citeonline{brooke1996}, com dez afirmações que alternam entre positivas e negativas. Essa alternância faz parte do instrumento e não deve ser corrigida. Para calcular o escore, subtrai-se 1 do valor marcado nos itens ímpares, subtrai-se de 5 o valor marcado nos itens pares, somam-se os resultados e multiplica-se o total por 2,5. O resultado vai de 0 a 100 e não é uma porcentagem. Nos estudos reunidos por \citeonline{sauro2012}, a média dos escores fica perto de 68 pontos. Para interpretar o escore, \citeonline{bangor2009} propõem faixas de aceitabilidade: abaixo de 50 pontos o resultado é considerado inaceitável, de 50 a 70 fica na faixa marginal e acima de 70 é considerado aceitável. A impressão geral usa a versão curta do questionário de \citeonline{schrepp2017}, com oito pares de adjetivos opostos em sete pontos; subtrai-se 4 do valor bruto e separam-se a dimensão pragmática e a hedônica. Como no instrumento original, o termo negativo de cada par fica sempre à esquerda e o positivo à direita, e os quatro itens pragmáticos vêm antes dos quatro hedônicos \cite{schrepp2023handbook}. A orientação fixa é própria da versão curta: na versão completa, de 26 itens, metade dos pares começa pelo termo positivo. A utilidade e a intenção de uso vêm do modelo de \citeonline{davis1989}, adaptado ao público de professores.

O bloco de riscos funciona ao contrário dos outros: nele, concordar significa apontar um problema. As afirmações são:

\begin{itemize}
    \item aprender em uma sintaxe personalizada pode dificultar a migração para uma linguagem de mercado;
    \item a liberdade de personalizar pode desviar a atenção do raciocínio algorítmico;
    \item o tempo gasto configurando a linguagem não compensa o ganho obtido.
\end{itemize}
 Os escores desse bloco são relatados separadamente e não podem ser somados aos do bloco de potencial pedagógico, porque as duas medidas apontam em sentidos opostos e uma anularia a outra. A mesma regra vale para a afirmação sobre sintaxes diferentes na mesma turma, que está no bloco de uso em sala. O bloco foi incluído de propósito: se o questionário trouxesse apenas afirmações favoráveis, a aprovação obtida diria muito pouco.

\subsection{Participantes, procedimento e questões éticas}

Os participantes do painel foram convidados diretamente e não são identificados. O formulário registra titulação, área de atuação, ferramentas de ensino que a pessoa já conhece e tempo de experiência com programação. Em uma avaliação feita por especialistas, esses dados fazem parte do resultado e são apresentados junto com os escores. Dois itens pedem que o respondente marque o que fez na plataforma e o que chegou a personalizar. Com isso é possível saber se uma avaliação negativa veio de quem usou o assistente até o fim ou de quem nem chegou a abri-lo.

Sobre o número de participantes: no modelo de \citeonline{nielsen1993}, cada novo participante revela uma parte menor de problemas ainda não encontrados, e por isso um grupo pequeno, de cinco a oito pessoas, já revela a maior parte dos problemas de usabilidade e basta para analisar as respostas abertas. Para apresentar o escore SUS com intervalo de confiança são necessários vinte respondentes ou mais. O painel teve oito, e por isso o resultado é apresentado como indicativo, sem afirmar que existe diferença entre subgrupos. Esse critério estava fixado antes da coleta. O painel é também uma amostra de conveniência, formada por convite.

Os dois formulários foram aplicados como pesquisa de opinião com participantes não identificados e não passaram pelo comitê de ética da instituição. O formulário dos especialistas abre com um termo de consentimento, apresentado como texto de leitura, sem item de marcação, e por isso as respostas desse formulário não trazem registro individual de consentimento. No formulário dos estudantes, o termo de assentimento é um item de resposta obrigatória, e as dezoito respostas trazem a opção ``Li e quero participar''. A sessão com estudantes envolveu menores de idade: nove tinham até 15 anos e nove tinham 16 ou 17. O formulário registra o assentimento do aluno, mas não a autorização de um responsável. Essas lacunas no registro do consentimento entram como limitação do estudo no Capítulo~\ref{cap:consideracoes}.

\subsection{Leitura dos resultados}
\label{subsec:leitura-resultados}

Os números apresentados a seguir vêm do resumo das respostas, gerado por um \textit{script} a partir das planilhas de respostas (Apêndice~\ref{ap:respostas}). Os blocos de análise, a ordem e os critérios de leitura foram definidos antes da coleta e foram mantidos; defini-los depois de ver os dados permitiria escolher o recorte mais favorável à plataforma. Os critérios de cálculo dos escores estão na Subseção~\ref{subsec:escalas-escores}.

As convenções de leitura valem para todas as subseções de resultados. Com oito respondentes, o texto usa contagens, médias e medianas, sem porcentagens. Nas afirmações em escala de 1 a 5, as respostas ``não usei'', ``não sei'' e ``não se aplica'' ficam fora da média e são contadas à parte. Nenhuma resposta foi excluída. Quando é preciso apontar um respondente, ele é indicado por R1 a R8, na ordem de chegada das respostas.

\subsection{Composição do painel}
\label{sec:resultado-painel}

Esta subseção usa as questões de perfil (\ref{qE:perfil-ini} a \ref{qE:perfil-fim}) e de uso da plataforma (\ref{qE:uso-ini} a \ref{qE:uso-fim}) do Apêndice~\ref{ap:formulario-especialistas}. Os oito respondentes foram convidados diretamente e não são identificados. Pela maior titulação declarada, três têm doutorado (a alternativa do formulário era ``doutorado ou superior''), dois têm graduação e três têm ensino médio, sem curso superior concluído. Cinco atuam como profissionais de desenvolvimento de software e três são professores do ensino superior, dois dos quais também lecionam no ensino técnico. Quatro lecionam ou já lecionaram disciplina que envolva programação: os três professores e um dos profissionais de desenvolvimento. Dois trabalham com programação há menos de 2 anos, três de 2 a 5 anos, dois de 6 a 10 anos e um há mais de 10 anos.

Sobre as ferramentas de ensino já conhecidas, cinco citaram o Portugol Studio ou o Portugol Webstudio, quatro citaram Scratch, Blockly ou similar, três citaram Beecrowd, URI ou outro juiz \textit{online}, dois citaram o VisuAlg, dois citaram Thonny, Python Tutor ou outro visualizador de execução, e um não conhecia nenhuma delas. A maior parte do painel, portanto, conhecia ao menos uma das ferramentas comparadas no Capítulo~\ref{cap:trabalhos-relacionados}.

O tempo de uso declarado foi curto para metade do painel: quatro usaram a plataforma por menos de 15 minutos, um de 15 a 30 minutos, um de 30 a 60 minutos e dois por mais de 1 hora. A Tabela~\ref{tab:resultado-atividades} mostra o que os respondentes disseram ter feito e o que personalizaram. Cinco criaram ou editaram uma linguagem no assistente e cinco escreveram e executaram um programa. Os recursos de turma foram os menos exercitados, junto com a importação de linguagens da comunidade.

\begin{table}[htbp]
\centering
\caption{Atividades declaradas pelos respondentes do painel e elementos da linguagem que personalizaram (respostas múltiplas, oito respondentes).}
\label{tab:resultado-atividades}
\footnotesize
\renewcommand{\arraystretch}{1.2}
\begin{tabular}{|p{10.5cm}|c|}
\hline
\textbf{O que fez na plataforma} & \textbf{Respondentes} \\\hline
Criou ou editou uma linguagem no assistente & 5 \\\hline
Escreveu e executou um programa & 5 \\\hline
Viu o código intermediário & 4 \\\hline
Usou a execução passo a passo & 4 \\\hline
Resolveu e submeteu um exercício & 3 \\\hline
Viu a tabela de \textit{tokens} & 2 \\\hline
Entrou em uma turma ou criou uma turma & 2 \\\hline
Importou uma linguagem da comunidade & 1 \\\hline
\textbf{O que personalizou na linguagem} & \textbf{Respondentes} \\\hline
Palavras dos tipos & 6 \\\hline
Modo de tipagem & 5 \\\hline
Entrada e saída & 4 \\\hline
Operadores em palavra & 3 \\\hline
Controle de fluxo & 2 \\\hline
Literais booleanos & 2 \\\hline
Delimitadores de bloco ou indentação & 2 \\\hline
Terminador de instrução & 1 \\\hline
Usou uma linguagem predefinida, sem alterar & 2 \\\hline
Não personalizou nada & 2 \\\hline
\end{tabular}
\legend{Fonte: dados da pesquisa (2026).}
\end{table}

Três características das respostas condicionam a leitura de tudo o que vem a seguir. A primeira é o pouco contato de parte do painel. R7 não marcou nenhuma atividade nem elemento personalizado, declarou menos de 15 minutos de uso e deu as notas mais baixas em vários blocos. R8 marcou apenas os recursos de turma e de exercício e escreveu, nas questões abertas, que não conseguiu explorar a plataforma. Os dois lecionam, e três dos quatro que lecionam estão entre os que usaram a plataforma por menos de 15 minutos. Parte das notas, portanto, reflete uma primeira impressão.

A segunda é o padrão de R2, que marcou o mesmo valor em todos os itens de cada bloco em escala: 3 no bloco de uso em sala, no bloco de recursos e no SUS, 6 nos oito pares do UEQ-S e 4 nos blocos de potencial pedagógico, de riscos, de qualidade técnica e de utilidade. Como o bloco de riscos tem sentido invertido, o mesmo 4 que concorda com os benefícios concorda também com os problemas. As respostas de R2 foram mantidas, e o efeito delas é indicado onde pesa.

A terceira é a relação entre o que foi avaliado e o que foi usado. Em vários itens o número de notas é maior que o número de pessoas que declararam ter usado o recurso. Um dos dois respondentes que marcaram não ter personalizado nada marcou também três elementos personalizados, e a resposta foi mantida como veio.

\subsection{Assistente de personalização da linguagem}
\label{sec:resultado-assistente}

Esta subseção trata das questões de avaliação dos recursos (\ref{qE:recursos-ini} a \ref{qE:recursos-fim}) e dos itens de qualidade técnica sobre o assistente (\ref{qE:qual-ini} a \ref{qE:qual-fim}) do Apêndice~\ref{ap:formulario-especialistas}. A Tabela~\ref{tab:resultado-recursos} traz os dez itens do bloco de avaliação dos recursos. Os três primeiros tratam do assistente. Nenhum item do bloco recebeu nota abaixo de 3.

\begin{table}[htbp]
\centering
\caption{Bloco de avaliação dos recursos (escala de 1 a 5; ``não usei'' contado à parte).}
\label{tab:resultado-recursos}
\footnotesize
\renewcommand{\arraystretch}{1.2}
\begin{tabular}{|p{7.9cm}|c|c|c|c|c|}
\hline
\textbf{Afirmação} & \textbf{n} & \textbf{Média} & \textbf{Mediana} & \textbf{Mín.--máx.} & \textbf{Não usei} \\\hline
Criar ou ajustar uma linguagem no assistente foi simples & 7 & 4,00 & 4 & 3--5 & 1 \\\hline
No assistente, ficou claro em que etapa eu estava & 7 & 4,43 & 5 & 3--5 & 1 \\\hline
O exemplo de código do assistente ajudou a entender o efeito de cada alteração & 7 & 4,57 & 5 & 3--5 & 1 \\\hline
Escrever programas com a linguagem configurada foi confortável & 8 & 4,12 & 4,5 & 3--5 & 0 \\\hline
As mensagens de erro ajudaram a identificar o que corrigir & 8 & 4,25 & 4,5 & 3--5 & 0 \\\hline
A tabela de \textit{tokens} é útil para entender como a plataforma lê o código & 6 & 4,50 & 5 & 3--5 & 2 \\\hline
O código intermediário é útil para entender o que o programa faz & 7 & 4,14 & 4 & 3--5 & 1 \\\hline
A execução passo a passo é útil para acompanhar o programa & 7 & 4,57 & 5 & 3--5 & 1 \\\hline
Resolver e submeter exercícios foi simples & 7 & 4,29 & 5 & 3--5 & 1 \\\hline
Importar uma linguagem da comunidade foi simples & 5 & 4,00 & 4 & 3--5 & 3 \\\hline
\end{tabular}
\legend{Fonte: dados da pesquisa (2026).}
\end{table}

Cada um dos três itens corresponde a uma decisão de interface descrita na Seção~\ref{sec:customizacao-dinamica}. A divisão da personalização em etapas é avaliada pelo item sobre a simplicidade de criar ou ajustar uma linguagem, que teve média 4,00: três respondentes marcaram 5, um marcou 4 e três marcaram 3. A barra de progresso com o título das etapas é avaliada pelo item sobre saber em que etapa se está, com média 4,43: cinco marcaram 5 e dois marcaram 3. A pré-visualização do código é avaliada pelo item sobre o exemplo, com média 4,57: cinco marcaram 5, um marcou 4 e um marcou 3. Nenhuma das três decisões foi contestada pelo painel. A mais bem avaliada foi a pré-visualização, e a que deixou mais gente no ponto neutro foi a simplicidade do fluxo como um todo. Os três que marcaram 3 nesse item são R2, com o padrão descrito na Seção~\ref{sec:resultado-painel}, e dois respondentes com menos de 15 minutos de uso, de modo que o dado não permite separar dificuldade real de pouco contato.

O ``não usei'' dos três itens veio do mesmo respondente, R8, que não abriu o assistente. Essa contagem mede quantas pessoas chegaram ao recurso, e não a qualidade dele. Sete pessoas deram nota aos três itens, mas só cinco declararam ter criado ou editado uma linguagem no assistente, de modo que duas notas vieram sem a declaração de uso correspondente.

A Tabela~\ref{tab:resultado-atividades} mostra quais partes do assistente foram mais exercitadas. As palavras dos tipos, o modo de tipagem e a entrada e saída, que ficam nas primeiras etapas, foram personalizados por seis, cinco e quatro respondentes. Os operadores, o controle de fluxo, os delimitadores de bloco e o terminador, que ficam nas etapas seguintes, foram personalizados por três, dois, dois e um. O formulário não permite saber se essas etapas foram percorridas sem alteração ou se não foram alcançadas.

Duas afirmações do bloco de qualidade técnica também tratam do assistente. A que diz que ele é compreensível para quem não tem formação em compiladores teve média 4,38, com mediana 4,5 e mínimo 3, entre os oito respondentes. A que diz que a validação prévia impede configurações de linguagem inválidas teve média 4,67 entre seis respondentes, com notas 4 e 5; os outros dois marcaram ``não sei''. A importação de linguagens da comunidade foi o recurso menos usado do bloco: três marcaram ``não usei'', cinco deram nota, com média 4,00, e só um declarou ter importado uma linguagem.

Os itens em escala avaliam o que acontece dentro do assistente. Uma resposta aberta, discutida na Seção~\ref{sec:resultado-abertas}, aponta um problema anterior a ele: a dificuldade de encontrar o botão que o abre.

\subsection{Uso do ambiente virtual de aprendizagem}
\label{sec:resultado-ava}

Esta subseção trata do bloco de uso em sala, respondido só por quem leciona (questões \ref{qE:sala-ini} a \ref{qE:sala-fim} do Apêndice~\ref{ap:formulario-especialistas}). O bloco de uso em sala foi respondido pelos quatro respondentes que lecionam ou já lecionaram. A Tabela~\ref{tab:resultado-ava} traz as nove afirmações. Com três ou quatro notas por linha, cada resposta pesa muito na média, e por isso as notas individuais são citadas no texto.

\begin{table}[htbp]
\centering
\caption{Bloco de uso em sala, respondido por quem leciona (escala de 1 a 5; ``não se aplica'' contado à parte).}
\label{tab:resultado-ava}
\footnotesize
\renewcommand{\arraystretch}{1.2}
\begin{tabular}{|p{7.8cm}|c|c|c|c|p{1.4cm}|}
\hline
\textbf{Afirmação} & \textbf{n} & \textbf{Média} & \textbf{Mediana} & \textbf{Mín.--máx.} & \textbf{Não se aplica} \\\hline
Criar uma turma e receber alunos pelo código de acesso foi simples & 3 & 3,67 & 3 & 3--5 & \centering\arraybackslash 1 \\\hline
Criar um exercício com casos de teste foi simples & 4 & 4,00 & 4 & 3--5 & \centering\arraybackslash 0 \\\hline
Travar uma linguagem específica em um exercício é útil para planejar as aulas & 4 & 4,50 & 5 & 3--5 & \centering\arraybackslash 0 \\\hline
A correção automática pelos casos de teste é útil para acompanhar a turma & 4 & 4,50 & 5 & 3--5 & \centering\arraybackslash 0 \\\hline
Eu usaria esta plataforma em uma disciplina introdutória de programação & 4 & 3,00 & 3 & 2--4 & \centering\arraybackslash 0 \\\hline
Eu conseguiria preparar uma aula com a plataforma sem apoio técnico externo & 4 & 4,00 & 4 & 3--5 & \centering\arraybackslash 0 \\\hline
O esforço de preparar a linguagem antes da aula compensa o ganho em sala & 3 & 3,00 & 3 & 2--4 & \centering\arraybackslash 1 \\\hline
Alunos de uma mesma turma usando sintaxes diferentes dificultaria o trabalho do professor (item de risco) & 3 & 4,00 & 4 & 3--5 & \centering\arraybackslash 1 \\\hline
O subconjunto da linguagem é suficiente para uma disciplina introdutória & 3 & 4,33 & 5 & 3--5 & \centering\arraybackslash 1 \\\hline
\end{tabular}
\legend{Fonte: dados da pesquisa (2026).}
\end{table}

Os recursos de turma e de correção foram bem avaliados. Travar uma linguagem em um exercício e a correção automática pelos casos de teste tiveram média 4,50, com três notas 5 e uma nota 3 em cada item. Criar um exercício com casos de teste teve média 4,00, com duas notas 5 e duas notas 3, e preparar uma aula sem apoio técnico externo também teve 4,00, com notas 3, 4, 4 e 5. Criar uma turma teve 3,67, com notas 3, 3 e 5. O subconjunto da linguagem foi considerado suficiente para uma disciplina introdutória por dois dos três que responderam, com nota 5; o terceiro marcou 3.

A intenção de adotar a plataforma ficou abaixo disso. A afirmação ``eu usaria esta plataforma em uma disciplina introdutória de programação'' teve média 3,00: um respondente concordou, com nota 4, dois ficaram no ponto neutro e um discordou, com nota 2. Ninguém marcou 5. A afirmação sobre o esforço de preparar a linguagem compensar o ganho em sala também teve média 3,00, com notas 2, 3 e 4. Os professores do painel, portanto, reconhecem utilidade nos recursos de turma e de correção, mas não declaram que adotariam a plataforma.

A pergunta sobre o momento da disciplina em que a plataforma seria mais útil aponta na mesma direção. Dois respondentes escolheram ``em paralelo à linguagem oficial, como apoio'', um escolheu ``antes da linguagem oficial da disciplina, como preparação'' e um escolheu ``não vejo espaço para ela na disciplina''. Quem escolheu essa última alternativa foi R7, que também deu a nota 2 para a adoção e para o esforço de preparação. O mesmo respondente deu nota 5 para o travamento de linguagem, para a correção automática, para a criação de exercícios e para a suficiência do subconjunto, e nota 5 no item de risco sobre sintaxes diferentes na mesma turma. A objeção de R7 não recai sobre os recursos de turma, avaliados com notas altas pelo mesmo respondente. Ela recai sobre o lugar de uma linguagem personalizável dentro da disciplina, e a implementação atual não tem resposta para isso.

Duas ressalvas pesam nesse bloco. Três dos quatro respondentes usaram a plataforma por menos de 15 minutos, e o quarto é R2, que marcou 3 em todas as linhas. Além disso, no painel inteiro só dois respondentes declararam ter entrado em uma turma ou criado uma, e um deles não leciona. As notas sobre criar turma e criar exercício vêm, em parte, de quem não percorreu esses fluxos. No bloco de recursos, que todos responderam, resolver e submeter exercícios teve média 4,29 entre sete respondentes, e três declararam ter resolvido e submetido um exercício.

\subsection{Usabilidade (SUS) e experiência de uso (UEQ-S)}
\label{sec:resultado-usabilidade}

Esta subseção trata do SUS (questões \ref{qE:sus-ini} a \ref{qE:sus-fim}) e do UEQ-S (questões \ref{qE:ueq-ini} a \ref{qE:ueq-fim}) do Apêndice~\ref{ap:formulario-especialistas}. O escore SUS médio do painel foi 75,6, com mediana 76,25, desvio padrão 16,5, mínimo 50,0 e máximo 95,0. A média fica acima da média de referência de 68 pontos \cite{sauro2012} e acima de 70, valor a partir do qual \citeonline{bangor2009} tratam o resultado como aceitável. Cinco respondentes ficaram acima de 68 e três ficaram abaixo. Com oito respondentes, o resultado é indicativo: não se calcula intervalo de confiança e não se comparam subgrupos, conforme o critério fixado no protocolo. A ressalva sobre a redação em português dos itens, registrada na Seção~\ref{subsec:limites-interpretacao}, também vale para essa comparação.

O escore 50,0 é o de R2, que marcou 3 nos dez itens. Sem essa resposta a média seria 79,3. O valor é informado para que o leitor veja o peso desse padrão; a média relatada continua sendo a dos oito.

Entre os itens positivos do SUS, o de menor concordância foi ``acho que gostaria de usar esta plataforma com frequência'', com média 3,50 e mediana 3. Em seguida vêm ``imagino que a maioria das pessoas aprenderia a usar esta plataforma rapidamente'', com 3,62, e ``achei a plataforma fácil de usar'', com 3,75. O de maior concordância foi ``considerei que as várias funções da plataforma estavam bem integradas'', com 4,25. Os cinco itens negativos tiveram médias entre 1,62 e 1,88, e nenhum respondente marcou mais que 3 em qualquer um deles. O item sobre uso frequente pede cuidado neste painel, porque os respondentes são professores e profissionais e a plataforma é voltada a iniciantes.

No UEQ-S, a dimensão pragmática teve média 1,78, a hedônica 1,88 e o valor geral 1,83, na faixa de $-3$ a $+3$. Na média do painel as duas dimensões ficaram próximas, com diferença de 0,10.

Em respondentes isolados as dimensões se afastam: em um deles a pragmática ficou bem abaixo da hedônica, em outro aconteceu o contrário, e um único respondente teve valor negativo na dimensão pragmática. Os dois pares pragmáticos com menor média bruta foram justamente complicada ou fácil e confusa ou clara, ambos com 5,62 em 7; atrapalhou ou apoiou teve 6,00 e ineficiente ou eficiente teve 5,88. Como este trabalho trata da interface, é nesses dois pares que está o sinal mais fraco, e ele vai na mesma direção do item do SUS sobre facilidade de uso.

Os valores do UEQ-S pedem cautela. Com oito respondentes, as médias servem só como indicação, e R2 marcou 6 nos oito pares, o que dá 2,00 nas duas dimensões.

\subsection{Potencial pedagógico e intenção de uso}
\label{sec:resultado-pedagogico}

Esta subseção trata do bloco de potencial pedagógico (questões \ref{qE:ped-ini} a \ref{qE:ped-fim}) e do bloco de utilidade e intenção de uso, adaptado do TAM (questões \ref{qE:tam-ini} a \ref{qE:tam-fim}), do Apêndice~\ref{ap:formulario-especialistas}. A Tabela~\ref{tab:resultado-pedagogico} reúne o bloco de potencial pedagógico e o bloco de utilidade e intenção de uso. Os dois foram respondidos pelos oito participantes. O que eles registram é a opinião de professores e profissionais sobre o que tende a acontecer com um iniciante, e não uma medida feita com iniciantes.

\begin{table}[htbp]
\centering
\caption{Blocos de potencial pedagógico e de utilidade e intenção de uso (escala de 1 a 5; ``não sei'' contado à parte).}
\label{tab:resultado-pedagogico}
\footnotesize
\renewcommand{\arraystretch}{1.2}
\begin{tabular}{|p{8.1cm}|c|c|c|c|c|}
\hline
\textbf{Afirmação} & \textbf{n} & \textbf{Média} & \textbf{Mediana} & \textbf{Mín.--máx.} & \textbf{Não sei} \\\hline
\multicolumn{6}{|l|}{\textbf{Potencial pedagógico}} \\\hline
Permitir que o aluno defina as palavras-chave tende a reduzir a barreira sintática no primeiro contato & 8 & 3,88 & 4 & 1--5 & 0 \\\hline
Usar palavras em português no lugar de termos em inglês facilita a compreensão do que o código faz & 8 & 4,38 & 5 & 2--5 & 0 \\\hline
A personalização tende a reduzir a frequência de erros de sintaxe cometidos por iniciantes & 7 & 4,14 & 5 & 1--5 & 1 \\\hline
Relatar os erros no vocabulário escolhido pelo próprio aluno facilita a correção do código & 8 & 4,25 & 5 & 2--5 & 0 \\\hline
Expor a tabela de \textit{tokens} tem valor didático para quem está aprendendo & 7 & 4,29 & 5 & 2--5 & 1 \\\hline
Expor o código intermediário tem valor didático para quem está aprendendo & 7 & 4,43 & 5 & 2--5 & 1 \\\hline
A execução passo a passo tem valor didático para quem está aprendendo & 8 & 4,75 & 5 & 4--5 & 0 \\\hline
O conjunto de elementos configuráveis é adequado ao propósito da ferramenta & 7 & 4,57 & 5 & 4--5 & 1 \\\hline
\multicolumn{6}{|l|}{\textbf{Utilidade e intenção de uso}} \\\hline
Esta plataforma é útil para quem está começando a programar & 8 & 4,38 & 5 & 2--5 & 0 \\\hline
Esta plataforma ajudaria alunos iniciantes a aprender lógica de programação mais rápido & 8 & 3,88 & 4,5 & 1--5 & 0 \\\hline
Comparada às ferramentas que eu conheço, esta plataforma oferece algo que elas não oferecem & 8 & 4,62 & 5 & 4--5 & 0 \\\hline
Recomendaria esta plataforma a um colega da área & 8 & 4,25 & 4,5 & 3--5 & 0 \\\hline
Vejo a plataforma como complemento, e não como substituto, à linguagem oficial da disciplina & 8 & 4,38 & 5 & 2--5 & 0 \\\hline
Indicaria esta plataforma a alguém que está começando a programar & 8 & 3,88 & 4,5 & 1--5 & 0 \\\hline
\end{tabular}
\legend{Fonte: dados da pesquisa (2026).}
\end{table}

As três primeiras afirmações do bloco pedagógico correspondem às premissas da introdução. Na primeira, sobre a redução da barreira sintática, seis respondentes concordaram, um ficou no ponto neutro e um discordou totalmente; a média foi 3,88. Na segunda, sobre o português facilitar a compreensão do código, sete concordaram e um discordou, com média 4,38. Na terceira, sobre a redução dos erros de sintaxe, seis concordaram, um discordou totalmente e um marcou ``não sei'', com média 4,14 entre sete. A maioria do painel concorda com as três premissas. A concordância não é unânime: R7 discordou das três, com notas 1, 2 e 1. A premissa central, a da barreira sintática, foi a de menor média entre os oito itens do bloco.

Os outros cinco itens ficaram acima de 4. O de maior média foi o valor didático da execução passo a passo, com 4,75 e nenhuma nota abaixo de 4. O conjunto de elementos configuráveis foi considerado adequado ao propósito da ferramenta, com média 4,57 e notas 4 e 5 entre os sete que responderam.

No bloco de utilidade e intenção de uso, a afirmação mais aceita foi a de que a plataforma oferece algo que as ferramentas conhecidas pelo respondente não oferecem, com média 4,62 e notas 4 e 5. As duas de menor média, ambas com 3,88, são as que mais se aproximam de uma promessa de resultado: que a plataforma ajudaria iniciantes a aprender lógica mais rápido e que o respondente a indicaria a alguém que está começando. Em cada uma, cinco concordaram, dois ficaram no ponto neutro e um discordou totalmente. O painel vê a plataforma como complemento à linguagem oficial da disciplina, e não como substituto dela, com média 4,38. Essa posição combina com a resposta ``em paralelo à linguagem oficial'' da Seção~\ref{sec:resultado-ava}.

Há uma diferença entre a intenção de uso declarada pelo painel inteiro e a declarada por quem leciona. Recomendar a plataforma a um colega teve média 4,25 entre os oito. Usá-la em uma disciplina introdutória teve média 3,00 entre os quatro que lecionam. São perguntas diferentes, feitas a grupos diferentes, e não cabe comparação estatística entre elas, mas a recomendação genérica foi mais alta que a disposição de adotar.

\subsection{Riscos e objeções levantados pelo painel}
\label{sec:resultado-riscos}

Esta subseção trata do bloco de riscos (questões \ref{qE:riscos-ini} a \ref{qE:riscos-fim} do Apêndice~\ref{ap:formulario-especialistas}) e da afirmação de risco do bloco de uso em sala (questão \ref{qE:sala-risco}). Nos itens desta seção, concordar significa apontar um problema. Eles são relatados separadamente e não entram em nenhuma média junto com o bloco de potencial pedagógico. A Tabela~\ref{tab:resultado-riscos} traz as três afirmações do bloco de riscos e a afirmação de risco do bloco de uso em sala, respondida só por quem leciona.

\begin{table}[htbp]
\centering
\caption{Itens de risco (escala de 1 a 5, em que concordar aponta um problema).}
\label{tab:resultado-riscos}
\footnotesize
\renewcommand{\arraystretch}{1.2}
\begin{tabular}{|p{5.8cm}|c|c|c|c|c|c|}
\hline
\textbf{Afirmação} & \textbf{n} & \textbf{Média} & \textbf{Mediana} & \textbf{Discordam} & \textbf{Neutros} & \textbf{Concordam} \\\hline
Aprender em uma sintaxe personalizada pode dificultar a migração posterior para uma linguagem de mercado & 8 & 3,00 & 3,5 & 3 & 1 & 4 \\\hline
A liberdade de personalizar pode desviar a atenção do aluno do raciocínio algorítmico & 7 & 2,57 & 2 & 4 & 0 & 3 \\\hline
O tempo gasto configurando a linguagem não compensa o ganho obtido & 7 & 1,86 & 2 & 6 & 0 & 1 \\\hline
Alunos de uma mesma turma usando sintaxes diferentes dificultaria o trabalho do professor (só quem leciona) & 3 & 4,00 & 4 & 0 & 1 & 2 \\\hline
\end{tabular}
\legend{Fonte: dados da pesquisa (2026).}
\end{table}

O risco de migração dividiu o painel. Quatro respondentes concordaram que aprender em uma sintaxe personalizada pode dificultar a passagem para uma linguagem de mercado, um ficou no ponto neutro e três discordaram. A implementação atual não responde a essa objeção. O que a plataforma tem são linguagens predefinidas inspiradas em linguagens conhecidas, descritas na Tabela~\ref{tab:linguagens-predefinidas}, que o professor pode travar em um exercício. Não há nenhum dado sobre a migração, e o risco continua em aberto. Só um estudo que acompanhe alunos depois da troca de linguagem poderia dizer algo sobre ele.

O risco de a personalização desviar a atenção do raciocínio algorítmico teve média 2,57, abaixo do ponto médio da escala, mas três dos sete que responderam concordaram, um deles totalmente. A implementação responde em parte. Quando o exercício ou a lista tem uma linguagem obrigatória, o seletor fica bloqueado e o aluno não pode trocar de linguagem durante a atividade, como descrito na Seção~\ref{sec:sistema-correcao}. Fora das atividades com linguagem travada, a escolha continua livre e a objeção permanece.

O risco de o tempo de configuração não compensar foi o menos aceito: seis discordaram e um concordou, R2, com média 1,86. As quatro linguagens predefinidas e a importação de linguagens da comunidade permitem começar sem passar pelo assistente. Esse resultado tem um contraponto na Seção~\ref{sec:resultado-ava}: entre os que lecionam, a afirmação de que o esforço de preparar a linguagem antes da aula compensa o ganho em sala ficou em 3,00. O painel em geral não vê o tempo de configuração como problema, mas os professores não chegam a dizer que o preparo compensa.

A afirmação sobre alunos da mesma turma com sintaxes diferentes foi respondida por três professores, com notas 3, 4 e 5. Nenhum discordou. Essa é a objeção a que a implementação responde de forma mais direta: o travamento de linguagem por exercício e por lista faz com que toda a turma resolva a atividade na mesma sintaxe, e os mesmos professores deram média 4,50 à utilidade desse travamento. Nas atividades em modo livre o problema continua. Nelas, a página de correção abre a submissão com a linguagem do aluno, como descrito na Subseção~\ref{subsec:correcao-professor}, o que permite ler o código como ele foi escrito, mas não elimina o trabalho de lidar com várias sintaxes.

\subsection{Respostas abertas}
\label{sec:resultado-abertas}

Esta subseção trata das questões abertas \ref{qE:abertas-ini} a \ref{qE:abertas-fim} do Apêndice~\ref{ap:formulario-especialistas}. O formulário tem quatro questões abertas e um item de autorização de citação. Sete respondentes autorizaram a citação sem identificação e um não autorizou. Cinco escreveram em pelo menos uma das questões, todos entre os que autorizaram. As respostas são poucas e curtas, e são relatadas aqui uma a uma.

Sobre o que funcionou melhor, houve quatro respostas. Duas apontam o assistente: ``configuração da linguagem'' e ``criação da linguagem''. Uma aponta os recursos de turma: ``criação de lista e acesso à lista criada''. A quarta fala da ``facilidade de utilização'' e da ``possibilidade de visualizar e compreender o funcionamento dos comandos de maneira dinâmica''.

Sobre os pontos de dificuldade, também houve quatro respostas, e duas delas interessam diretamente à navegação. Um respondente escreveu: ``dificuldade em visualizar o botão de criar linguagem''. Na questão sobre o que mudaria, o mesmo respondente pediu um ``botão de criar linguagem mais óbvio/chamativo''. Esse respondente é R4, que deu notas 4, 5 e 5 aos três itens do assistente e a nota 2, a mais baixa do painel, ao item do SUS ``achei a plataforma fácil de usar''. A combinação indica que, para esse respondente, o problema esteve em chegar ao assistente e não em usá-lo. O relato não diz em qual tela a dificuldade ocorreu. Na versão atual do código, a página inicial, onde fica a IDE, tem duas entradas para o assistente. A primeira está no título acima da IDE: só as palavras ``Crie sua própria'' formam o botão, com fundo discreto, e o restante da frase, ``linguagem de programação'', é texto comum. A segunda é um ícone de adição sem rótulo visível, colocado no painel lateral de linguagem depois da lista de lexemas, dentro da área de rolagem, de modo que pode ficar fora da vista até que o usuário role o painel; o mesmo ícone aparece abaixo da lista do seletor de linguagens. No ícone, o nome ``Criar linguagem'' só aparece na dica que surge quando o ponteiro passa sobre ele. Um botão com rótulo escrito, ``Nova linguagem'', existe na página de linguagens, fora da IDE.

O outro relato é de R8: ``não consegui explorar significativamente a plataforma''. O mesmo respondente escreveu, nas outras duas questões, que não explorou as mensagens de erro e que não explorou o suficiente para sugerir alterações. A resposta não diz o que impediu a exploração, se a falta de tempo ou a dificuldade de se orientar na interface. R8 declarou menos de 15 minutos de uso, marcou ``não usei'' em seis dos dez recursos e não abriu o assistente. Mesmo sem a causa, o relato mostra que um professor, em uma primeira visita curta, chegou às turmas e aos exercícios e não chegou à personalização da linguagem, que é a funcionalidade central da plataforma.

Uma terceira resposta sobre dificuldade trata da linguagem, e não da interface: segundo o respondente, ``pode haver alguma dificuldade para compreender a sintaxe e a estrutura dos comandos, especialmente para quem já está habituado a outras linguagens de programação''. A quarta diz apenas que não houve dificuldade.

Sobre as mensagens de erro, um respondente sugeriu ``mensagens de erro mais descritivas, indicando com clareza onde está o problema e, sempre que possível, sugerindo como corrigi-lo''. Os outros dois que responderam disseram que nenhuma mensagem foi confusa ou que não exploraram essa parte. A sugestão convive com notas altas nos itens em escala: 4,25 para as mensagens de erro terem ajudado a identificar o que corrigir e 4,71, entre sete respondentes, para a clareza delas no bloco de qualidade técnica.

Sobre o que mudariam, além do pedido de um botão mais visível, um respondente propôs ampliar os recursos interativos, ``com mais exemplos práticos, exercícios progressivos e desafios''.

